% ----------------------------------------------------------
% Introdução
% ----------------------------------------------------------
\chapter{Introdução}
\label{cap:introducao}
% ----------------------------------------------------------

Desde a primeira olimpíada moderna em Atenas 1896, mais de 137 mil atletas competiram por medalhas olímpicas ao longo de 125 anos de história (até Tokyo 2020, realizado em 2021). Por trás desses resultados individuais, existe uma rede complexa de relações competitivas que revela padrões latentes de dominância, rivalidades e agrupamentos estruturais não capturados por métricas tradicionais. A modelagem e a análise de redes complexas têm sido amplamente aplicadas para compreender sistemas complexos em diversas áreas do conhecimento, revelando padrões estruturais e dinâmicas não capturados por análises tradicionais baseadas em métricas isoladas \cite{Ahmed2023ComplexNetworks}. No contexto esportivo, esta abordagem oferece perspectivas complementares ao permitir a modelagem de competições como sistemas relacionais, onde atletas e suas interações competitivas podem ser representados como grafos direcionados e ponderados \cite{Radicchi2016OlympicDominance}.

A análise de redes complexas revela propriedades estruturais complementares não capturadas por métricas tradicionais. Enquanto estatísticas convencionais avaliam atletas isoladamente através de tempo, pontos ou recordes, a análise de redes captura hierarquias de influência, padrões de dominância competitiva e agrupamentos estruturais \cite{Wasche2017EmergingParadigm, LopezFelip2018ClusterPhase}.

Esta metodologia permite compreender o ecossistema competitivo através de interações entre atletas de diferentes modalidades e países ao longo do tempo, considerando competições esportivas como sistemas dinâmicos complexos \cite{Davids2013}

Trabalhos recentes têm aplicado redes complexas ao contexto esportivo \cite{Motegi2012NetworkBased, Shiue2025ApplyingPageRank}. Contudo, a maioria concentra-se em uma análise isolada de modalidades específicas, sem explorar diferenças estruturais sistemáticas entre tipos de interdependência competitiva ao longo de períodos históricos extensos.

O presente trabalho preenche esta lacuna aplicando ferramentas de ciência de redes a dados olímpicos históricos (1896-2021), desenvolvendo uma metodologia que integra a modelagem de grafos competitivos, a análise de centralidade e a detecção de comunidades. Esta abordagem visa revelar padrões estruturais através de métricas avançadas como \textit{PageRank} e algoritmo de \textit{Louvain} em seis modalidades representativas que abrangem três tipos estruturais: esportes mistos (natação e atletismo), esportes coletivos (basquetebol e futebol) e esportes individuais puros (judô e boxe).

Estabelecido o contexto e as lacunas da literatura científica existente, este trabalho busca responder à seguinte questão de pesquisa:

\section{Questão de Pesquisa}

Este trabalho investiga a seguinte questão:

\textit{Como a estrutura de redes competitivas olímpicas difere entre esportes com diferentes tipos de interdependência competitiva (individual, coletivo, misto)?}

Especificamente, busca-se identificar que padrões de dominância e agrupamento estrutural emergem ao longo de mais de um século de competições olímpicas.

\section{Objetivos}

Este trabalho objetiva aplicar a análise de redes complexas a dados históricos olímpicos (1896-2021) para identificar e quantificar diferenças estruturais entre modalidades com tipos distintos de interdependência competitiva (individual, coletivo, misto) através de métricas de centralidade, detecção de comunidades e caracterização multidimensional de agrupamentos estruturais.

Os objetivos específicos compreendem:

\begin{itemize}
    \item Construir 16 redes competitivas direcionadas e ponderadas a partir de 9.350 medalhistas olímpicos, segregando por esporte e gênero nas modalidades de natação, atletismo, basquetebol, futebol, judô e boxe;
    \item Aplicar métricas de centralidade (\textit{PageRank}, \textit{degree}, \textit{betweenness}, \textit{closeness}) para identificar atletas estruturalmente importantes nas redes competitivas;
    \item Detectar e caracterizar comunidades nas redes através do algoritmo de \textit{Louvain}, analisando suas propriedades temporais, geográficas e de performance;
    \item Implementar \textit{dashboard} interativo como artefato complementar para exploração e visualização dos resultados;
    \item Analisar comparativamente as propriedades estruturais das redes entre esportes individuais, coletivos e mistos.
\end{itemize}

\section{Contribuições}

Este trabalho apresenta três contribuições principais para a literatura de análise de redes aplicada ao esporte:

\textbf{Caracterização Multidimensional de Comunidades:} A análise de comunidades desenvolvida integra métricas estruturais de tamanho e modularidade com a caracterização temporal através de entropia de Shannon e heterogeneidade de performance através do coeficiente de variação de \textit{PageRank}. Esta sistematização permite compreender não apenas a existência de agrupamentos estruturais, mas suas características distintivas no contexto competitivo de cada modalidade.

\textbf{Análise Comparativa Multi-Esporte:} A análise estrutural de competições com diferentes tipos de interdependência (individual, coletivo, misto) ao longo de 125 anos de história olímpica permite a identificação de padrões sistemáticos relacionados ao formato competitivo.

\textbf{\textit{Dashboard} Interativo para Exploração:} A implementação de uma interface web interativa oferece um artefato complementar que demonstra a viabilidade de ferramentas acessíveis para exploração de redes complexas por audiências não-técnicas. O \textit{dashboard} integra a análise de centralidades, a visualização de comunidades, a análise temporal e comparações entre esportes, facilitando a identificação de padrões e a validação de hipóteses sobre estruturas competitivas.

\section{Organização do Trabalho}

Este trabalho está organizado em cinco capítulos. O Capítulo 2 apresenta a revisão bibliográfica, abordando os fundamentos teóricos de redes complexas, a modelagem de redes competitivas esportivas, as métricas de centralidade e os algoritmos de detecção de comunidades, além de trabalhos relacionados que aplicaram a análise de redes ao contexto esportivo. O Capítulo 3 descreve a metodologia adotada, detalhando o processo de coleta e preparação dos dados olímpicos, a construção das redes competitivas, os critérios de ponderação de arestas, e os procedimentos de análise através de métricas de centralidade e detecção de comunidades. O Capítulo 4 apresenta os resultados obtidos, incluindo a caracterização estrutural das redes, os rankings de atletas por \textit{PageRank}, a análise de comunidades detectadas e as comparações entre modalidades esportivas. O Capítulo 5 conclui o trabalho, sintetizando as principais contribuições, limitações e direções para pesquisas futuras.